\section{The exotic-hadron model in detail}\label{app:exotics}

This appendix gives the model, the fit, the full spectra, the threshold catalogue and the details of the H9 tests behind
Section~\ref{sec:exotics}. Headline values are not repeated here; Table~\ref{tab:exo-a} and Table~\ref{tab:exo-mint} are their home.

\subsection{The colour-spin Hamiltonian in the H8 trees}\label{app:exotics-model}

\paragraph{Colour and spin.} In Eq.~\eqref{eq:exo-H}, $\boldsymbol\lambda_i\!\cdot\!\boldsymbol\lambda_j=\sum_{a=1}^8\lambda^a_i
\lambda^a_j$ is the colour product of the Gell-Mann matrices, with $-\lambda^{a*}$ for an antiquark, and $\boldsymbol\sigma$ are the Pauli
matrices. The state is the H8 tree. Every internal node below the root couples its two branches to a colour (anti)triplet, with
quadratic Casimir $4/3$; the root couples to a singlet; every node carries a spin. The colour factors that follow are
$\langle\boldsymbol\lambda_i\!\cdot\!\boldsymbol\lambda_j\rangle=-16/3$ for a quark and antiquark in a singlet and $-8/3$ for a pair in
a diquark or a baryon. A quark and an antiquark in the two halves of $[qq][\bar q\bar q]$ have $-4/3$, the two diquarks of
$([qq][qq])\bar q$ have $-2/3$ per pair, and its antiquark has $-4/3$ with each quark \chk{exotics:E}. Within a fixed colour tree, node
spins that give the same total $J$ are mixed and the Hamiltonian is diagonalised. For neutral states with hidden flavour, the
eigenvectors are sorted by $C$ parity, the exchange of the two halves. Two identical $s$, $c$ or $b$ quarks in a diquark have spin 1
(Pauli), and a light $[ud]$ pair has isospin equal to its spin. The S-wave pentaquark multiplets have the multiplicities of Maiani, Polosa
and Riquer~\cite{maiani2015}.

\paragraph{The trees span the colour space.} The products $\mathbf3\otimes\mathbf3=\mathbf{\bar3}\oplus\mathbf6$,
$\mathbf{\bar3}\otimes\mathbf{\bar3}=\mathbf3\oplus\mathbf{\bar6}$ and $\mathbf3\otimes\mathbf{\bar3}=\mathbf1\oplus\mathbf8$ each contain
the conjugate triplet or the singlet. So any quarks and antiquarks that can form a singlet can be coupled pairwise until two conjugate
halves meet at the root \tagP. For a tetraquark the two branches of H8 exhaust the colour space. The singlets of $qq\bar q\bar q$ form a
two-dimensional space, and the compact tree $[q_1q_2]_{\bar3}[\bar q_3\bar q_4]_3$ and the molecular tree $(q_1\bar q_3)_1(q_2\bar
q_4)_1$ are two vectors of it with $|\langle\text{compact}|\text{molecular}\rangle|^2=1/3$ \chk{exotics:E}. Every tetraquark colour state
is therefore a superposition of the two trees, and H8 cannot fail on a single state. No numerology is used: the five quarks of a
pentaquark are not connected with the five parts of the minting, since the axioms supply no map between them.

\paragraph{Couplings.} Mesons and baryons have separate effective masses $m^m_f$ and $m^b_f$ ($f=q,s,c,b$, $q=u,d$). The coefficient of
$\boldsymbol\sigma_i\!\cdot\!\boldsymbol\sigma_j$ is $v^m_{ij}=2b\,k_{ij}/(m^m_im^m_j)$ in a meson and $v^b_{ij}=a\,k_{ij}/(m^b_im^b_j)$ in
a baryon, so that
\begin{equation}\label{eq:exo-C}
C_{ij}=\tfrac{3}{16}\,v^m_{ij}\ \ (\text{quark--antiquark pair}),\qquad C_{ij}=\tfrac38\,v^b_{ij}\ \ (\text{quark--quark pair})
\end{equation}
returns $v^m$ and $v^b$ for the colour factors $-16/3$ and $-8/3$. The factor $k_{ij}$ is 1 except for $(c,s)$ and $(b,s)$, where it is
KR's enhancement $k_s=a_{cs}/a$. KR apply it in baryons; we share it with mesons, where KR found the same enhancement of the
$D_s^*-D_s$ splitting. Heavy--heavy pairs take $v^m_{c\bar c}=[M(J/\psi)-M(\eta_c)]/4$ and $v^m_{b\bar b}=[M(\Upsilon)-M(\eta_b)]/4$.
$v^m_{b\bar c}$ is interpolated between them with $|\psi(0)|^2\propto\mu^p$ in the reduced mass $\mu$~\cite[Sec.~VII]{karliner2014}, and
a heavy--heavy quark pair has the same $C$. The binding terms are
\begin{equation}\label{eq:exo-B}
B_{ij}=B^{(1)}_{ij}\,\frac{\langle\boldsymbol\lambda_i\!\cdot\!\boldsymbol\lambda_j\rangle}{-16/3},
\end{equation}
with colour-singlet values $B^{(1)}$ for $(c,\bar s)$, $(b,\bar s)$, $(c,\bar c)$, $(b,\bar b)$ and $(b,\bar c)$ fitted to mesons, and none
for pairs with a $u$ or $d$ quark. In a baryon or diquark this is KR's rule $B(QQ)=B(Q\bar Q)/2$. For a pair in different halves of a
multiquark it gives $B^{(1)}/4$: KR's stated assumption, a potential that factorises into a colour factor and one spatial factor,
applied to every pair.

\subsection{Variants of the mass convention}\label{app:exotics-variants}

Two variants were fixed with the model and are reported for every state, never scored. \emph{P-in} drops the binding between constituents
of different halves. \emph{J} replaces the two mass sets by one universal set and adds a constant $S$ per string junction
\cite{karliner2017,karliner2017a}: none for a meson, one for a baryon, two for a tetraquark and three for a pentaquark, the internal
nodes below the root. Variant J fits ordinary hadrons as well as the primary model ($\sigma_{\rm model}=8.2$\,MeV, $S=160.6$\,MeV). It
shifts the tetraquarks by 90 to 113\,MeV and the lowest pentaquarks by 204 to 210\,MeV, ten to twenty-five times $\sigma_{\rm model}$.
The ordinary hadrons cannot choose between the conventions, and this choice dominates every compact mass. Of the nine scored
confrontations, the primary model holds 4, P-in 2 and J 2 \chk{exotics:A}.

\subsection{The fit and the hosting check}\label{app:exotics-fit}

\paragraph{Inputs.} The 42 hadrons are $\pi,\rho,\omega,K,K^*,\phi$, $D^{(*)},D_s^{(*)},B^{(*)},B_s^{(*)},\eta_c,J/\psi,\eta_b,\Upsilon(1S),B_c$,
the light octet and decuplet baryons, and the singly charmed and bottom baryons $\Lambda_Q,\Sigma_Q^{(*)},\Xi_Q,\Xi_Q',\Xi_Q^*,
\Omega_Q^{(*)}$ (no $\Omega_b^*$). The masses are the values tabulated by KR~\cite{karliner2014} or current values of the Particle Data
Group~\cite{navas2024}, read state by state; the $\Xi_b'$ uses LHCb's mass difference~\cite{lhcb2015h}. Isospin averages are taken, since
the model has no isospin breaking; $\eta$ and $\eta'$ are left out because of flavour mixing, as by KR.

\paragraph{Uncertainty.} $\sigma_{\rm model}=\sqrt{\mathrm{RSS}/(N-p)}$ with $N=42$ and $p=18$. A prediction carries
$\sigma_{\rm pred}^2=\sigma_{\rm model}^2+g^{\mathsf T}\Sigma g$, with $g$ the gradient of the predicted mass,
$\Sigma=\sigma_{\rm model}^2(\mathsf D^{\mathsf T}\mathsf D)^{-1}$ the parameter covariance and $\mathsf D$ the Jacobian of the fitted
masses. The residuals have root-mean-square 5.5\,MeV for mesons and 6.6\,MeV for baryons, and the largest is $-16.3$\,MeV, for the
$\Sigma$ (Fig.~\ref{fig:exo-fit}). Fitting with each hadron left out in turn and predicting it gives 11.2\,MeV over 37 hadrons; the five
that fix a parameter exactly are skipped. Our parameters (Table~\ref{tab:exo-par}) differ from KR's mainly in $m^b_c$ and
$B^{(1)}_{c\bar c}$, because KR fix each parameter on one or two hadrons and we fit all at once with current charmonium masses
\chk{exotics:F}.

\paragraph{Hosting check.} With KR's published parameters the code reproduces KR's tables to their quoted precision: 0.15\,MeV for the
heavy hadrons, and the rounding to 1\,MeV for the light ones. It gives $\Xi_{cc}$ at 3626.7\,MeV, and $T(bb\bar u\bar d)=10389.3$,
$T(cc\bar u\bar d)=3882.15$ and $T(bc\bar u\bar d)=7133.75$\,MeV against KR's 10389.4, 3882.2 and 7133.7\,MeV~\cite{karliner2017}. It
also reproduces the fully charmed $0^{++}$ estimate of Karliner, Nussinov and Rosner, 6191.5\,MeV~\cite{karliner2017a}, from their
formula \chk{exotics:H}.

\begin{table}[tbp]\centering\small
\caption{Fitted parameters (1$\sigma$ from the covariance scaled by $\sigma_{\rm model}$) and KR's values~\cite{karliner2014}, all in
MeV; the rows $b$ and $a$ give $b/(310\,\text{MeV})^2$ and $a/(363\,\text{MeV})^2$ \chk{exotics:F}.}
\label{tab:exo-par}
\begin{tabular}{@{}lrrlrr@{}}
\toprule
parameter & fit & KR & parameter & fit & KR \\
\midrule
$m^m_q$ & $309.6\pm2.1$ & 310 & $m^b_q$ & $364.0\pm1.4$ & 363 \\
$m^m_s$ & $481.8\pm3.1$ & 483 & $m^b_s$ & $536.0\pm1.8$ & 538 \\
$m^m_c$ & $1658.1\pm5.2$ & 1663.3 & $m^b_c$ & $1715.9\pm4.4$ & 1710.5 \\
$m^m_b$ & $5001.6\pm5.4$ & 5003.8 & $m^b_b$ & $5047.8\pm4.5$ & 5043.5 \\
$b$ & $79.3\pm1.5$ & 80 & $a$ & $51.1\pm1.0$ & 50 \\
$k_s$ & $1.79\pm0.12$ & 1.65 (baryons) & $B^{(1)}_{c\bar s}$ & $-66.1\pm6.7$ & $-69.9$ \\
$B^{(1)}_{b\bar s}$ & $-80.4\pm7.0$ & $-83.6$ & $B^{(1)}_{c\bar c}$ & $-247.5\pm12.3$ & $-258.0$ \\
$B^{(1)}_{b\bar b}$ & $-558.4\pm12.6$ & $-562.8$ & $B^{(1)}_{b\bar c}$ & $-334.6\pm12.4$ & $-341.5$ \\
$v^m_{c\bar c}$ & $28.2\pm2.9$ & 28.3 & $v^m_{b\bar b}$ & $15.4\pm2.9$ & 15.6 \\
\bottomrule
\end{tabular}
\end{table}

\begin{figure}[tbp]\centering
\includegraphics[width=0.49\linewidth]{exotic_fit.pdf}
\caption{Fit of the colour-spin model to 42 ordinary ground-state hadrons: model minus measured mass against mass, mesons (circles) and
baryons (squares); open circles are the residuals when the hadron is left out of the fit; the band is $\pm\sigma_{\rm model}=\pm8.1$\,MeV.
The diamond is the held-out $\Xi_{cc}^{++}$ with $\pm\sigma_{\rm pred}$. Data: KR's tables~\cite{karliner2014} and the Particle Data
Group~\cite{navas2024}. What to look for: no trend with mass, and the held-out baryon within $2\sigma$ \chk{exotics:F}.}
\label{fig:exo-fit}
\end{figure}

\subsection{The full spectra}\label{app:exotics-spectra}

Table~\ref{tab:exo-full} adds to Table~\ref{tab:exo-a} the same trees with KR's own parameters, the two variants, the states that are not
scored and the isospin partners; Fig.~\ref{fig:exo-levels} shows the pentaquark and four-charm levels \chk{exotics:A}.

\begin{table}[tbp]\centering\footnotesize\setlength{\tabcolsep}{2.8pt}
\caption{All compact trees: the model (KR's convention, our fit) with $\sigma_{\rm pred}$, the same tree with KR's parameters, the
variants P-in and J, the datum, the pull, and model mass minus the lowest S-wave two-hadron threshold. Data as in Table~\ref{tab:exo-a},
$X(3872)$~\cite{lhcb2020k} and $Z_c(3900)$~\cite{besiii2013h}. Masses in MeV \chk{exotics:A}.}
\label{tab:exo-full}
\begin{tabular}{@{}llrrrrrrr@{}}
\toprule
state & tree, $J^{P(C)}$ & model & KR par. & P-in & J & data & $z$ & $M-T$ \\
\midrule
$T_{cc}^+$ & $[cc]_1[\bar u\bar d]_0$, $1^+$ & $3897\pm12$ & 3882 & 3897 & 3995 & 3874.74 & $+1.86$ & $+22$ \\
$X_0(2900)$ & $[ud]_0[\bar c\bar s]_0$, $0^+$ & $2755\pm10$ & 2754 & 2755 & 2861 & 2866 & $-9.28$ & $+392$ \\
$T_{c\bar s0}(2900)$ & $[cq][\bar s\bar q]$, $0^+$ & $2809\pm9$ & 2812 & 2825 & 2916 & 2908 & $-4.06$ & $+701$ \\
$X(6900)$ & $[cc]_1[\bar c\bar c]_1$, $2^{++}$ & $6425\pm29$ & 6383 & 6672 & 6515 & 6905 & $-15.2$ & $+231$ \\
$P_c(4312)$ & lowest $([cq][qq])\bar c$ & $4277\pm11$ & 4264 & 4339 & 4485 & 4311.9 & $-3.06$ & $+354$ \\
$P_c(4440)^\dagger$ & $([cq][qq]_1)\bar c$, $\frac12^-$ & $4428\pm11$ & 4446 & 4396 & 4485 & 4440.3 & $-0.98$ & $+506$ \\
$P_c(4457)^\dagger$ & $([cq][qq]_1)\bar c$, $\frac12^-$ & $4462\pm11$ & 4456 & 4490 & 4485 & 4457.3 & $+0.38$ & $+539$ \\
$P_{\psi s}(4338)$ & $([cs][ud]_0)\bar c$, $\frac12^-$ & $4392\pm11$ & 4382 & 4470 & 4602 & 4338.2 & $+4.81$ & $+292$ \\
$P_{cs}(4459)^\dagger$ & $([cs][ud]_0)\bar c$, $\frac32^-$ & $4460\pm11$ & 4442 & 4470 & 4602 & 4458.8 & $+0.13$ & $+361$ \\
\midrule
\multicolumn{9}{@{}l}{\emph{not scored}}\\
$X(6900)$ & $[cc]_1[\bar c\bar c]_1$, $0^{++}$ & $6340\pm26$ & 6298 & 6588 & 6430 & 6905 & $-19.1$ & $+372$ \\
$X(3872)$ & $[cq][\bar c\bar q]$, $1^{++}$ & $4108\pm11$ & 4093 & 4170 & 4212 & 3871.70 & $+21.0$ & $+237$ \\
$Z_c(3900)$ & $[cq][\bar c\bar q]$, $1^{+-}$ & $3985\pm11$ & 3969 & 4047 & 4087 & 3899 & $+6.72$ & $+749$ \\
\midrule
\multicolumn{9}{@{}l}{\emph{live}}\\
$T_{bb}^-$ & $[bb]_1[\bar u\bar d]_0$, $1^+$ & $10399\pm12$ & 10389 & 10399 & 10512 & -- & -- & $-205$ \\
$T_{bc}$ & $[bc]_0[\bar u\bar d]_0$, $0^+$ & $7146\pm11$ & 7134 & 7146 & 7251 & -- & -- & $+2$ \\
$T_{bc}$ & $[bc]_1[\bar u\bar d]_0$, $1^+$ & $7180\pm12$ & 7168 & 7180 & 7285 & -- & -- & $-9$ \\
$P_{cc}$ & $([cc]_1[ud]_0)\bar q$, $\frac12^-$ & $4232\pm12$ & 4215 & 4232 & 4436 & -- & -- & $+471$ \\
$cc\bar c\bar c$ & $[cc]_1[\bar c\bar c]_1$, $1^{+-}$ & $6369\pm27$ & 6326 & 6616 & 6459 & -- & -- & $+288$ \\
$T_{cc}$, $I=1$ & $[cc]_1[\bar q\bar q]_1$, $0^+,1^+,2^+$ & 4042, 4071, 4131 & & & & -- & -- & $+312,+200,+117$ \\
$T_{bb}$, $I=1$ & $[bb]_1[\bar q\bar q]_1$, $0^+,1^+,2^+$ & 10583, 10593, 10613 & & & & -- & -- & $+25,-11,-37$ \\
\bottomrule
\end{tabular}
\end{table}

\paragraph{Doubly heavy states.} $T_{cc}^+$ is LHCb's pole $360\pm40$\,keV below the $D^{*+}D^0$ threshold of 3875.10\,MeV, with
$J^P=1^+$ and isospin zero~\cite{lhcb2022h,lhcb2022i}. For $T_{bb}^-$, lattice QCD finds binding energies relative to $B\bar B^*$ of
$189\pm10$~\cite{francis2017}, $128\pm26$~\cite{leskovec2019}, $83\pm22$~\cite{aoki2023}, $112\pm13$~\cite{hudspith2023} and
$116^{+36}_{-30}$\,MeV~\cite{tripathy2025}; against our 205\,MeV the pulls run from $-1.0$ to $-5.1$. These are calculations and enter no
verdict. For the $1^+$ $bc\bar u\bar d$ state, lattice QCD finds binding of 15 to 61\,MeV~\cite{francis2019} or $\Delta E=-43^{+15}_{-25}$\,MeV
\cite{padmanath2024}, pulls $+1.1$ and $+1.8$ against our $-9$\,MeV. Alexandrou et al.\ find shallow bound or virtual states a few MeV or
less below both $B^{(*)}\bar D$ thresholds~\cite{alexandrou2024}.

\paragraph{Open charm and strangeness.} $X_0(2900)$ is a spin-0 $\bar cd\bar su$ state in $B^+\to D^+D^-K^+$~\cite{lhcb2020i}, and
$T_{c\bar s0}(2900)^{++,0}$ an isospin triplet with $J^P=0^+$ in $D_s^+\pi^\pm$~\cite{lhcb2023h}. The spin-1 $X_1(2900)$ has negative parity,
needs an orbital excitation and lies outside the S-wave model.

\paragraph{Four charm quarks.} The S-wave levels lie 146 to 231\,MeV above $J/\psi J/\psi$. LHCb's $X(6900)$ is $6905\pm11\pm7$\,MeV
without interference with the continuum~\cite{lhcb2020h}. CMS's no-interference fit finds structures at $6552\pm16$, $6927\pm10$ and
$7287^{+21}_{-19}$\,MeV~\cite{cms2024h}; ATLAS's model~A finds $6.41^{+0.11}_{-0.09}$, $6.63^{+0.09}_{-0.05}$ and
$6.86^{+0.03}_{-0.04}$\,GeV~\cite{atlas2023h}; LHCb sees a broad enhancement from 6.2 to 6.8\,GeV. Only the lowest ATLAS structure lies
near the S-wave levels; this was noticed after the comparison and is not a test.

\paragraph{Pentaquarks.} The $([cq][qq])\bar c$ multiplet has 10 levels with isospin $1/2$, from 4277 to 4584\,MeV. The lowest has a
spin-0 light diquark and $J=1/2$. Variant P-in, without the $c\bar c$ binding across the halves, gives $z=+2.0$ for $P_c(4312)$. For
$udsc\bar c$ with isospin 0 we compute the two trees $([cs][ud]_0)\bar c$ and $([cq][qs])\bar c$; the lowest $J=1/2$ level belongs to the
first. $P_{cs}(4459)$ lies 1.5\,MeV from a $J=3/2$ level.

\begin{figure}[tbp]\centering
\includegraphics[width=\linewidth]{exotic_levels.pdf}
\caption{Levels of the compact trees (model, KR's convention, our fit). Left: hidden-charm pentaquarks $([cq][qq])\bar c$ with $I=1/2$
by $J$ (all of negative parity), coloured by the spin of the light diquark, against LHCb's $P_c(4312)$, $P_c(4440)$,
$P_c(4457)$~\cite{lhcb2019h} (solid) and the $\Sigma_c\bar D^{(*)}$, $\Sigma_c^*\bar D$ thresholds (dashed). Centre: $udsc\bar c$ with
$I=0$ in the two trees $([cs][ud]_0)\bar c$ and $([cq][qs])\bar c$ against $P_{\psi s}(4338)$~\cite{lhcb2023j} and $P_{cs}(4459)$
\cite{lhcb2021h}, with the $\Xi_c\bar D^{(*)}$ thresholds. Right: $[cc]_1[\bar c\bar c]_1$ for the model and both variants against the
di-$J/\psi$ structures of LHCb (model I)~\cite{lhcb2020h}, CMS (no interference)~\cite{cms2024h} and ATLAS (model A)~\cite{atlas2023h}.
What to look for: the lowest model levels sit 35 to 55\,MeV from the lowest observed pentaquarks, the levels are dense enough above
4.4\,GeV to match any peak, and the S-wave four-charm levels lie far below 6.9\,GeV \chk{exotics:A}.}
\label{fig:exo-levels}
\end{figure}

\subsection{The threshold catalogue}\label{app:exotics-thresholds}

\paragraph{Catalogue.} The 26 states, each from its primary publication, are $\chi_{c1}(3872)$~\cite{lhcb2020k,lhcb2013h},
$Z_c(3900)$~\cite{besiii2013h,besiii2017h}, $Z_c(4020)$~\cite{besiii2013i}, $Z_c(4430)$~\cite{lhcb2014h}, $Z_{cs}(3985)$~\cite{besiii2021h},
$Z_{cs}(4000)$, $Z_{cs}(4220)$, $\chi_{c1}(4140)$, $\chi_{c1}(4274)$, $\chi_{c0}(4500)$, $\chi_{c0}(4700)$, $X(4685)$, $X(4630)$
\cite{lhcb2021i,lhcb2017h}, $X(3960)$~\cite{lhcb2023i}, $\psi(4230)$~\cite{besiii2017i}, $Z_b(10610)$, $Z_b(10650)$~\cite{belle2012h},
$T_{cc}^+$~\cite{lhcb2022i}, $X_0(2900)$, $X_1(2900)$~\cite{lhcb2020i}, $T_{c\bar s0}(2900)$~\cite{lhcb2023h}, $P_c(4312)$, $P_c(4440)$,
$P_c(4457)$~\cite{lhcb2019h}, $P_{\psi s}(4338)$~\cite{lhcb2023j} and $X(6900)$~\cite{lhcb2020h}.

\paragraph{Thresholds.} We take every pair of hadrons narrower than 60\,MeV, from 83 hadrons and their antiparticles~\cite{navas2024}:
light, open-charm and open-bottom mesons including the narrow P-wave ones, charmonia and bottomonia below open flavour, and light,
charmed and bottom baryons. A pair counts if its quark flavours add up to the state's and it can form the state in an S wave. That
requires $J$ in $|j_1-j_2|,\dots,j_1+j_2$ and parity $P_1P_2$. For neutral states with hidden flavour it also requires $C=C_1C_2$ for two
self-conjugate hadrons, or $(-1)^S$ for a particle and its own antiparticle. Unknown quantum numbers impose nothing; isospin and $G$ parity
are not imposed.

\paragraph{Null and statistic.} Each mass is placed uniformly in $[M-W,M+W]$ with $W=150$\,MeV, and
\begin{equation}\label{eq:exo-p}
p_i=\frac{1}{2W}\Big|\,[M_i-W,M_i+W]\cap\bigcup_k[T_k-d_i,T_k+d_i]\Big|
\end{equation}
is the probability that the nearest allowed threshold $T_k$ lies at most as far as the observed $d_i$. It is computed exactly and agrees
with 20000 random placements per state. The $p_i$ are combined by Fisher's method, $X^2=-2\sum_i\ln p_i\sim\chi^2_{52}$; $X^2=101.1$.
The windows contain 5.5 allowed thresholds on average; this look-elsewhere count is built into each $p_i$ \chk{exotics:B}.

\paragraph{Robustness.} Within 5\,MeV of an allowed threshold lie 7 states, where 3.8 are expected. Restricted to the 17 states with
measured $J^P$, $p=1.8\times10^{-4}$. Perturbing the 45 threshold masses we could not re-read by 1\,MeV (3\,MeV) moves $p$ within
$7.0\times10^{-6}$ to $7.5\times10^{-5}$ ($5.0\times10^{-7}$ to $1.4\times10^{-4}$). Twelve of the 26 states lie below their nearest
threshold and fourteen above.

\begin{table}[tbp]\centering\scriptsize\setlength{\tabcolsep}{2.6pt}
\caption{Threshold proximity: mass (MeV), $J^{P(C)}$ (? unknown), nearest allowed S-wave threshold, $M-T$ (MeV) and $p_i$ of
Eq.~\eqref{eq:exo-p} for $W=150$\,MeV \chk{exotics:B}.}
\label{tab:exo-b}
\begin{tabular}{@{}lrlrrr@{\hspace{8pt}}lrlrrr@{}}
\toprule
state & $M$ & $J^{P(C)}$ & threshold & $M-T$ & $p_i$ & state & $M$ & $J^{P(C)}$ & threshold & $M-T$ & $p_i$ \\
\midrule
$\chi_{c1}(3872)$ & 3871.7 & $1^{++}$ & $D^0\bar D^{*0}$ & $+0.0$ & 0.0001 & $\psi(4230)$ & 4222.0 & $1^{--}$ & $\omega\chi_{c0}$ & $+24.6$ & 0.59 \\
$Z_c(3900)$ & 3899.0 & $1^{+}$ & $\psi(3770)\pi$ & $-14.3$ & 0.39 & $Z_b(10610)$ & 10607.2 & ? & $B\bar B^*$ & $+2.8$ & 0.075 \\
$Z_c(4020)$ & 4022.9 & ? & $D^*\bar D^*$ & $+5.8$ & 0.11 & $Z_b(10650)$ & 10652.2 & ? & $B^*\bar B^*$ & $+2.7$ & 0.055 \\
$Z_c(4430)$ & 4475.0 & $1^{+}$ & $D_1\bar D_1$ & $-373.3$ & 0.72 & $T_{cc}^+$ & 3874.7 & $1^+$ & $D^0D^{*+}$ & $-0.4$ & 0.0072 \\
$Z_{cs}(3985)$ & 3982.5 & ? & $\bar D^0D_s^{*+}$ & $+5.5$ & 0.34 & $X_0(2900)$ & 2866.0 & $0^+$ & $D^{*-}K^{*+}$ & $-36.1$ & 0.24 \\
$Z_{cs}(4000)$ & 4003.0 & $1^{+}$ & $J/\psi K^{*+}$ & $+14.2$ & 0.33 & $X_1(2900)$ & 2904.0 & $1^-$ & $\bar D_1^0K^0$ & $-15.7$ & 0.11 \\
$Z_{cs}(4220)$ & 4216.0 & $1^{?}$ & $\psi(2S)K^+$ & $+36.2$ & 0.89 & $T_{c\bar s0}(2900)$ & 2908.0 & $0^+$ & $D^{*+}K^{*+}$ & $+5.9$ & 0.039 \\
$\chi_{c1}(4140)$ & 4118.0 & $1^{++}$ & $J/\psi\phi$ & $+1.6$ & 0.022 & $P_c(4312)$ & 4311.9 & ? & $\Sigma_c^+\bar D^0$ & $-5.6$ & 0.25 \\
$\chi_{c1}(4274)$ & 4294.0 & $1^{++}$ & $\omega\psi(2S)$ & $-174.8$ & 0.99 & $P_c(4440)$ & 4440.3 & ? & $p\chi_{c1}$ & $-8.6$ & 0.55 \\
$\chi_{c0}(4500)$ & 4474.0 & $0^{++}$ & $\omega\psi(2S)$ & $+5.2$ & 0.10 & $P_c(4457)$ & 4457.3 & ? & $\Lambda_c(2595)\bar D^0$ & $+0.2$ & 0.022 \\
$\chi_{c0}(4700)$ & 4694.0 & $0^{++}$ & $\phi\psi(2S)$ & $-11.6$ & 0.42 & $P_{\psi s}(4338)$ & 4338.2 & $\frac12^-$ & $\Xi_c^+D^-$ & $+0.8$ & 0.063 \\
$X(4685)$ & 4684.0 & $1^{++}$ & $\phi\psi(2S)$ & $-21.6$ & 0.49 & $X(6900)$ & 6905.0 & $?^{?+}$ & $\chi_{c0}\chi_{c1}$ & $-20.4$ & 0.86 \\
$X(4630)$ & 4626.0 & $1^{-+}$ & $D_s^*\bar D_{s1}(2536)$ & $-21.3$ & 0.79 & $X(3960)$ & 3956.0 & $0^{++}$ & $\eta'\eta_c$ & $+14.1$ & 0.33 \\
\bottomrule
\end{tabular}
\end{table}

\begin{figure}[tbp]\centering
\includegraphics[width=\linewidth]{exotic_thresholds.pdf}
\caption{Threshold proximity of the 26 exotic hadrons of Table~\ref{tab:exo-b}. Left: distance to the nearest allowed S-wave
two-hadron threshold (bars, logarithmic) and the median distance under the null of uniform placement within $\pm150$\,MeV (ticks).
Right: empirical distribution of the per-state $p_i$ against the uniform null, for $W=150$ and 300\,MeV, with Fisher's combined $p$. What
to look for: the excess of small $p_i$ comes from a handful of states, two of them within 1\,MeV of threshold \chk{exotics:B}.}
\label{fig:exo-thresholds}
\end{figure}

\subsection{The tests of H9 in detail}\label{app:exotics-mint}

\paragraph{I1.} $P_c(4380)$ was reported above 9 standard deviations in 2015~\cite{lhcb2015i}; its verification awaits an amplitude analysis
\cite{lhcb2019h}. $P_c(4337)$ is evidence at 3.1 to 3.7 standard deviations~\cite{lhcb2022j}, and $P_{cs}(4459)$ at 3.1~\cite{lhcb2021h}. In
$\Lambda_c^+D_s^-$, the production ratios of $P_{\psi s}(4338)$ and $P_{cs}(4459)$ are below 0.12 and 0.20 at 95\% CL~\cite{lhcb2025h}. The
hidden row of the $2\times2$ table holds 4 established and 3 unconfirmed claims, the open row 0 and 3; Fisher's exact $p$ is 0.17 one-sided
and 0.20 two-sided \chk{exotics:I1}.

\paragraph{I2.} Binding energies are taken relative to isospin-averaged thresholds: $P_c(4312)$ lies 8.8\,MeV below $\Sigma_c\bar D$,
$P_c(4440)$ and $P_c(4457)$ lie 21.7 and 4.7\,MeV below $\Sigma_c\bar D^*$, and $P_{\psi s}(4338)$ lies 1.9\,MeV above $\Xi_c\bar D$. The
threshold masses are KR's averages~\cite{karliner2014} and the Particle Data Group's $B^*$~\cite{navas2024}. In the variant, the binding
scales with the reduced mass of the pair (a ratio of 2.5 to 2.6), since a heavier pair binds more deeply in the same potential. The
compact multiplet $([bq][qq])\bar b$ spans 10883 to 11141\,MeV (variant J: lowest 11100\,MeV), and the lowest $udsb\bar b$ level with
$J=1/2$ lies at $10993\pm11$\,MeV \chk{exotics:I2}.

\begin{table}[tbp]\centering\footnotesize
\caption{Hidden-bottom pentaquarks as molecules: the $P_c$ and $P_{\psi s}$ binding carried over unchanged, or scaled with the reduced
mass. Masses in MeV \chk{exotics:I2}.}
\label{tab:exo-pb}
\begin{tabular}{@{}llrrrr@{}}
\toprule
partner of & bottom threshold & $T$ & binding & $M$ (same binding) & $M$ (scaled) \\
\midrule
$P_c(4312)$ & $\Sigma_b\bar B$ & 11093.7 & $+8.8$ & 11084.9 & 11070.8 \\
$P_c(4440)$ & $\Sigma_b\bar B^*$ & 11139.0 & $+21.7$ & 11117.4 & 11084.5 \\
$P_c(4457)$ & $\Sigma_b\bar B^*$ & 11139.0 & $+4.7$ & 11134.4 & 11127.3 \\
$P_{\psi s}(4338)$ & $\Xi_b\bar B$ & 11072.1 & $-1.9$ & 11074.0 & 11077.0 \\
\bottomrule
\end{tabular}
\end{table}

\paragraph{I4: eligibility.} Eligibility was fixed from quark content and $J^P$ only, and the list was written into the script before
any splitting was computed. A pair is two exotic states with the same quark content (isospin partners counted once), each reported as an
observation, with equal or unknown parity and spins that differ by at most one unit or are unknown. The 25 eligible pairs are: three
among $Z_c(3900)$, $Z_c(4020)$, $Z_c(4430)$; three among the $Z_{cs}$; fifteen among the six $c\bar cs\bar s$ states of positive parity
($X(3960)$, $\chi_{c1}(4140)$, $\chi_{c1}(4274)$, $\chi_{c0}(4500)$, $X(4685)$, $\chi_{c0}(4700)$); $Z_b(10610)/Z_b(10650)$; $P_c(4312)$
with each of $P_c(4440)$ and $P_c(4457)$; and the CMS pair at 6552 and 6927\,MeV~\cite{cms2024h}. Not eligible are $X(4630)$,
$X(3872)/\psi(4230)$ and $X_0/X_1(2900)$, whose parities differ. The two-peak fit of $P_{cs}(4459)$ is excluded because LHCb can neither
confirm nor refute it~\cite{lhcb2021h}, and $X(7100)$, $P_c(4337)$ and $P_c(4380)$ because they are not observations.

\paragraph{I4: test.} A pair passes if $|\Delta-\alpha M/2|<1.96\,\sigma_\Delta$, with $\alpha=1/137.035999177$~\cite{codata2022}. Under a
null in which $\Delta$ is uniform between 0 and 300\,MeV, a pair passes by chance with the probability that this interval covers; the sum
over the 25 pairs, 4.2, is the look-elsewhere count, and the Poisson probability of two or more passes is 0.92. As the QCD baseline we
take the colour-spin model's splitting of two S-wave levels with the same quantum numbers, where the model has such a pair
(Table~\ref{tab:exo-alpha}, Fig.~\ref{fig:exo-alpha}) \chk{exotics:I4}.

\begin{table}[tbp]\centering\footnotesize
\caption{The splitting rule $\Delta=\alpha M/2$: measured splitting, the rule, its pull, and the colour-spin model's splitting of the
corresponding pair of S-wave levels with its pull. Shown: the pairs that pass and the pairs with a model baseline; the other 18 eligible
pairs fail by 4.8 to 38 standard deviations (Fig.~\ref{fig:exo-alpha}). Below the line, two pairs that are not eligible. Masses in MeV
\chk{exotics:I4}.}
\label{tab:exo-alpha}
\begin{tabular}{@{}lrrrrr@{}}
\toprule
pair & $\Delta$ & $\alpha M/2$ & pull & model $\Delta$ & pull \\
\midrule
$Z_c(3900)/Z_c(4020)$ & $123.9\pm6.7$ & 14.5 & $+16.3$ & $131.6\pm8.9$ & $-0.7$ \\
$Z_b(10610)/Z_b(10650)$ & $45.0\pm2.5$ & 38.8 & $+2.5$ & $143.8\pm8.9$ & $-10.7$ \\
$Z_{cs}(3985)/Z_{cs}(4000)$ & $20.5\pm11.4$ & 14.6 & $+0.5$ & $63.9\pm8.3$ & $-3.1$ \\
$X(4685)/\chi_{c0}(4700)$ & $10.0\pm18.0$ & 17.1 & $-0.4$ & -- & -- \\
$\chi_{c0}(4500)/\chi_{c0}(4700)$ & $220.0\pm10.8$ & 16.7 & $+18.7$ & $32.5\pm8.4$ & $+13.7$ \\
$X(3960)/\chi_{c0}(4500)$ & $518.0\pm12.0$ & 15.4 & $+42.0$ & $32.5\pm8.4$ & $+33.2$ \\
$X(3960)/\chi_{c0}(4700)$ & $738.0\pm15.0$ & 15.8 & $+48.1$ & $32.5\pm8.4$ & $+41.0$ \\
\midrule
$P_c(4440)/P_c(4457)$, suggested the rule & 17.0 & 16.2 & -- & -- & -- \\
$P_{cs}(4459)$, unresolved two-peak fit & $12.9\pm4.6$ & 16.3 & -- & -- & -- \\
\bottomrule
\end{tabular}
\end{table}

\begin{figure}[tbp]\centering
\includegraphics[width=0.49\linewidth]{exotic_alpha.pdf}
\caption{Splittings of the 25 eligible pairs of exotic hadrons against their mean mass (logarithmic scale; lower error bars clipped at
the axis), with the rule $\alpha M/2$ (line), the colour-spin model's splitting where it has a corresponding pair of S-wave levels
(triangles), and the pair that suggested the rule (star, not counted). Data: Table~\ref{tab:exo-b} and CMS~\cite{cms2024h}. What to look
for: the eligible pairs scatter over two orders of magnitude, and the two within $1.96\sigma$ of the line have errors as large as the
prediction \chk{exotics:I4}.}
\label{fig:exo-alpha}
\end{figure}

\subsection{Caveats in full}\label{app:exotics-caveats}

\begin{itemize}[leftmargin=*,itemsep=1pt]
\item KR's mass convention dominates the compact masses (Appendix~\ref{app:exotics-variants}).
\item $\sigma_{\rm model}$ measures ordinary hadrons. Carrying the model to four and five quarks adds uncertainty it does not contain: the
size of a diquark, binding across the halves, and colour configurations outside the tree, such as $\mathbf6\otimes\mathbf{\bar6}$,
whose mixing would lower the lowest levels.
\item Only S-wave ground states are computed, so a mismatch refutes an S-wave assignment, not every compact reading.
\item The threshold test inherits the selection of its catalogue, and its look-elsewhere count depends on the threshold list.
\item Lattice results for $T_{bb}$ scatter by 100\,MeV.
\item The server of the Particle Data Group was not reachable. Its current values were read through a search of its listings, and 45 of
the 83 threshold hadrons were not re-read; perturbing them moves no conclusion.
\end{itemize}
